\documentclass[twocolumn,11pt,a4paper]{article}

% ======================================================================
%  A BEHAVIOURAL ISOMORPHISM CALCULUS FOR VIDEO EFFECT COMPOSITION
%  Backward Trajectory Completion, Catalytic Chains, and a Domain-Specific
%  Language for Media Effect Encoding
%
%  Self-contained. No reliance on unpublished work; all citations are to
%  standard, externally published literature.
% ======================================================================

\usepackage[utf8]{inputenc}
\usepackage[T1]{fontenc}
\usepackage{lmodern}
\usepackage[a4paper,margin=2cm]{geometry}
\usepackage{amsmath,amssymb,amsthm,mathtools}
\usepackage{bm}
\usepackage{graphicx}
\usepackage{enumitem}
\usepackage{booktabs}
\usepackage{array}
\usepackage{microtype}
\usepackage{xcolor}
\usepackage{listings}
\usepackage[colorlinks=true,linkcolor=blue!45!black,citecolor=blue!45!black,
            urlcolor=blue!45!black]{hyperref}
\usepackage[round]{natbib}

% ---- theorem environments ----
\newtheoremstyle{mee}{6pt}{6pt}{\itshape}{}{\bfseries}{.}{.5em}{}
\theoremstyle{mee}
\newtheorem{theorem}{Theorem}[section]
\newtheorem{lemma}[theorem]{Lemma}
\newtheorem{proposition}[theorem]{Proposition}
\newtheorem{corollary}[theorem]{Corollary}

\theoremstyle{definition}
\newtheorem{definition}[theorem]{Definition}
\newtheorem{axiom}{Axiom}
\newtheorem{principle}[theorem]{Principle}
\newtheorem{example}[theorem]{Example}

\theoremstyle{remark}
\newtheorem{remark}[theorem]{Remark}

% ---- DSL listing style ----
\lstdefinelanguage{MEE}{
  morekeywords={clip,scene,goal,brand,dispatch,when,do,play,acts,like,
                compose,render,emit,via,at,for,fps,aspect,duration},
  sensitive=true,
  morecomment=[l]{//},
  morestring=[b]",
}
\lstset{
  language=MEE,
  basicstyle=\small\ttfamily,
  keywordstyle=\bfseries,
  commentstyle=\itshape\color{gray!70},
  frame=none,
  xleftmargin=1em,
  columns=fullflexible,
  keepspaces=true,
}

% ---- macros ----
\newcommand{\Real}{\mathbb{R}}
\newcommand{\Nat}{\mathbb{N}}
\newcommand{\Out}{\mathcal{X}}
\newcommand{\State}{\mathcal{S}}
\newcommand{\BState}{\mathcal{B}}
\newcommand{\floor}{\beta}
\newcommand{\pot}{\kappa}
\newcommand{\Eff}{\varepsilon}
\newcommand{\Chain}{\Gamma}
\newcommand{\target}{x^{\dagger}}
\newcommand{\src}{x^{0}}
\newcommand{\ISO}{\phi}
\newcommand{\Desc}{\mathcal{D}}
\newcommand{\Phys}{\mathcal{P}}
\newcommand{\IR}{\mathsf{IR}}
\newcommand{\dist}{d}
\newcommand{\Sig}{\Sigma}
\newcommand{\compose}{\mathbin{\triangleright}}

\title{\bfseries A Behavioural Isomorphism Calculus\\
for Video Effect Composition:\\
\large Backward Trajectory Completion, Catalytic Chains,\\
and a Domain-Specific Language for Media Effect Encoding}

\author{%
  \normalsize Pugachev Cobra Project --- Working Paper\\
  \small \texttt{media-effect-encoder}
}
\date{\today}

\begin{document}
\maketitle

\begin{abstract}
\noindent
We develop a formal calculus for the composition of video effects from a
single structural primitive: the \emph{behavioural isomorphism}, a
sequence of incremental state transformations connecting a footage clip to
a target behaviour description. The user specifies two things---a source
clip and a description of how the output should \emph{behave}---and the
system derives the connecting chain. We prove, as consequences of a
covering-radius floor theorem, that such chains are always finite,
non-unique, and that intermediate states need not individually resemble
either endpoint. We formalise effects as \emph{information catalysts}
whose power composes multiplicatively, prove a saturation dichotomy
(Borel--Cantelli type) governing when a chain can reach the target cell,
and prove that coherence of a chain is decidable from ordinal data alone.
The result is a domain-specific language---the \emph{Media Effect
Encoder} (MEE)---in which a user describes physical behaviours
(\emph{acts like water}, \emph{acts like a drum skin}) rather than
parameter values, and the compiler derives and verifies the isomorphism
chain. We describe the language surface, the intermediate representation,
the type-checker as a static assessor of chain validity, and the two-phase
compile/render architecture. The theory is ordinal: it asserts structural
properties of chains, never absolute parameter magnitudes.
\end{abstract}

% ======================================================================
\section{Introduction}
\label{sec:intro}
% ======================================================================

\subsection{The problem}

Video post-production tools are parameterised. A colour grade is a matrix
of numbers; a blur is a radius; a distortion is an amplitude and a
frequency. Users navigate menus of named presets---``warm look,''
``desaturate,'' ``filmic grain''---each a frozen bag of parameter values
assembled by a tool author. Two deficiencies follow from this design.

First, named presets conflate two distinct levels of description: what a
scene should \emph{do} (its behavioural character) and how that behaviour
is realised as a specific parameter setting. A user who wants a shot to
``move like water'' does not know, and should not need to know, the exact
frequency, amplitude, and damping coefficient of the displacement map that
achieves it. The named-preset vocabulary forces them to navigate a library
assembled by someone else rather than describe their own intent.

Second, the vocabulary is closed. New effects require new library entries,
new documentation, and user training on new names. The space of possible
behaviours is open; the space of named presets is not.

We argue that both deficiencies share a root: the representation is at the
wrong level of abstraction. Physical behaviour is the correct level. A
user knows that water oscillates, refracts, and reflects; that a drum skin
exhibits radial propagation from a strike point, damped overtones, and
membrane tension. These are behavioural descriptions composed from physical
phenomena whose parameter instantiation is a compiler problem, not a user
problem. The user specifies the behaviour; the compiler derives the chain
of primitives that realises it.

\subsection{The core model}

A footage clip occupies a state in a high-dimensional perceptual space.
A behaviour description names a target region of that space---a cell of
states that share the described character. The \emph{effect chain} is an
isomorphism: a finite sequence of incremental transformations that moves
the clip's state from its initial position toward the target cell. Each
step is small---a single physical primitive---and the chain composes them
multiplicatively. The user never specifies the steps; they specify the
endpoints. The chain is derived by the compiler.

This model has an exact precedent in the theory of backward trajectory
completion \citep{poincare1890,wiener1948}: given a known endpoint, a
trajectory to it can be recovered by working backwards from the terminus.
The efficiency of backward completion is exponential in the depth of the
state hierarchy \citep{sutton1988,bellman1957}: a chain of $k$ steps
backward costs $O(k)$, whereas forward enumeration of all paths of length
$k$ costs $O(n^k)$ where $n$ is the branching factor. The compiler is a
backward-trajectory engine: it knows where the chain must end (the target
cell) and derives the steps.

A further consequence of the backward model is \emph{path opacity}: two
chains with the same endpoints but different intermediate steps are
indistinguishable by any test of the output \citep{kolmogorov1965}. This
is not a defect; it is a feature. The user cannot tell whether the
``water'' effect was realised as oscillation-then-refraction or as
refraction-then-oscillation because the outputs are equivalent. The
compiler is free to choose whichever ordering is cheapest to render.

\subsection{Contributions}

\begin{enumerate}[leftmargin=1.4em,itemsep=2pt]
\item \textbf{Behavioural isomorphism chains} (\S\ref{sec:iso}). We
  formalise the connection between a footage state and a behaviour
  description as a chain of physical primitives, and prove that such
  chains are always finite and non-unique.
\item \textbf{The covering floor and chain existence}
  (\S\ref{sec:floor}). We prove, from a covering-radius argument, that
  every target cell has a finite-length chain reaching its interior, and
  that no chain can reach its centre exactly (the floor theorem).
\item \textbf{Catalytic composition and saturation} (\S\ref{sec:power}).
  We prove the multiplicative power law for stacked primitives, the
  diminishing-returns corollary, and a Borel--Cantelli saturation
  dichotomy.
\item \textbf{Coherence of a chain} (\S\ref{sec:coherence}). We prove
  that a chain is coherent if and only if its support graph contains a
  directed cycle of length $\ge3$, and that coherence is decidable from
  ordinal data alone.
\item \textbf{The MEE language} (\S\ref{sec:language}). We define the
  syntax and semantics of the Media Effect Encoder DSL: source clips,
  behaviour descriptions, effect chains as pipelines, a goal block, and
  a dispatch table for multi-format output.
\item \textbf{The intermediate representation and compiler}
  (\S\ref{sec:compiler}). We define the IR and the two-phase
  compile/render architecture, with the type-checker as a static assessor
  of chain coherence and saturation before any frame is rendered.
\end{enumerate}

% ======================================================================
\section{State Spaces and Physical Primitives}
\label{sec:state}
% ======================================================================

\subsection{The media state space}

\begin{definition}[Media state space]
\label{def:state}
A \emph{media state space} is a triple $(\State, \dist, \Sig)$ where
$\State$ is a non-empty set of media states (audiovisual configurations
of a clip at a given moment), $\dist:\State\times\State\to[0,\Sig]$ is a
bounded pseudometric, and $\Sig\in(0,\infty)$ is the maximal perceptual
distance. We write $\src\in\State$ for the \emph{source state} (the
unprocessed footage) and $\BState\subseteq\State$ for a
\emph{behaviour cell}: a positive-diameter region whose members all
exhibit the same described behaviour.
\end{definition}

A behaviour cell is the formal correlate of a behavioural description. The
cell $\BState_{\text{water}}$ is the set of all media states that exhibit
the physical character of a water surface: oscillatory displacement,
refracted light, specular reflection. The description names the cell; the
cell is not a point.

\begin{axiom}[Finite primitive vocabulary]
\label{ax:finite}
The set of physical primitives available to the compiler is finite:
$|\Phys|<\infty$. Each primitive $p\in\Phys$ is a transformation
$p:\State\to\State$ with bounded displacement $\dist(\src, p(\src))\le
\delta_p$ for some $\delta_p>0$.
\end{axiom}

Axiom~\ref{ax:finite} is not a limitation; it is a design choice. The
set of physical phenomena that a rendering engine can instantiate is
finite (oscillation, refraction, reflection, damping, compression, shear,
diffusion, scatter, and their combinations). What is open is not the
vocabulary but the parameterisation: each primitive carries continuous
parameters (frequency, amplitude, axis, damping coefficient) that the
compiler instantiates.

\begin{definition}[Behavioural description]
\label{def:desc}
A \emph{behavioural description} $\Desc$ is a finite set of physical
phenomena $\{b_1,\dots,b_m\}\subseteq\Phys$ that jointly characterise a
target behaviour. It specifies the \emph{what} (the phenomena) without
specifying the \emph{how} (the parameters). The cell $\BState_\Desc$ is
the set of all states exhibiting all phenomena in $\Desc$.
\end{definition}

``Acts like water'' is the description
$\Desc_{\text{water}}=\{\text{oscillate},\text{refract},\text{reflect},
\text{ripple-propagate}\}$. Each element names a class of physical
transformations; the class is not a single parameter setting but an
equivalence class of settings that all produce the phenomenon. The cell
$\BState_{\text{water}}$ is therefore genuinely a region, not a point.

% ======================================================================
\section{Behavioural Isomorphism Chains}
\label{sec:iso}
% ======================================================================

\subsection{Effects as primitive applications}

\begin{definition}[Effect]
\label{def:effect}
An \emph{effect} is a pair $\Eff=(p,\theta)$ where $p\in\Phys$ is a
physical primitive and $\theta$ is a parameter assignment. The effect acts
on a state as $\Eff(x)=p_\theta(x)$ where $p_\theta$ is $p$ instantiated
at $\theta$.
\end{definition}

An effect is atomic: it does one physical thing with one parameter
setting. ``Oscillate vertically at 0.8 Hz with 12-pixel amplitude'' is an
effect; ``water'' is not. The user describes behaviours; the compiler
produces effects.

\begin{definition}[Isomorphism chain]
\label{def:chain}
A \emph{behavioural isomorphism chain} from source $\src$ toward cell
$\BState$ is a finite sequence of effects
$\Chain=(\Eff_1,\dots,\Eff_n)$ such that:
\begin{enumerate}[label=(\roman*),leftmargin=1.8em,itemsep=1pt]
\item Each step is small: $\dist(\Eff_i(x_{i-1}), x_{i-1})\le\delta_{p_i}$
  where $x_{i-1}$ is the state after $i-1$ steps.
\item The chain is directed: each step reduces $\dist(x_i, \BState)$
  weakly and at least one step reduces it strictly.
\item The chain terminates: $\dist(x_n, \BState)\le\floor$ where
  $\floor>0$ is the covering floor (Definition~\ref{def:floor}).
\end{enumerate}
The chain \emph{realises} the description $\Desc$ if every primitive in
$\Desc$ appears as the physical component of at least one $\Eff_i$.
\end{definition}

The realisation condition formalises the user's intent: a chain for
``water'' must include at least one oscillation effect, at least one
refraction effect, and so on. The chain need not apply them in any fixed
order; it need only apply each phenomenon at least once.

\begin{remark}[Steps need not individually resemble the target]
\label{rem:subgoal}
Condition (i) requires each step to be small relative to the current
state, not relative to the target cell. An intermediate state $x_i$ may
look nothing like a water surface; it is only the global output $x_n$
that must lie within $\floor$ of $\BState$. This is the formal content of
the walking analogy: a person walking to a destination passes through
states (one foot raised, weight shifted, arms swung) that individually
look nothing like \emph{arrived}. The trajectory's validity is a global
property of the chain, not a local property of each step.
\end{remark}

\subsection{Non-uniqueness and path opacity}

\begin{theorem}[Path non-uniqueness]
\label{thm:nonunique}
For any source $\src$ and behaviour cell $\BState$ with
$\dist(\src,\BState)>0$, there exist at least two distinct isomorphism
chains realising the same description $\Desc$.
\end{theorem}

\begin{proof}
Let $\Chain=(\Eff_1,\dots,\Eff_n)$ be a chain realising $\Desc$.
Permuting two adjacent effects $\Eff_i,\Eff_{i+1}$ that act on
independent perceptual channels (e.g., a displacement on the vertical
axis and a colour temperature shift) produces a chain with the same
endpoint set of primitives but different intermediate states. Since the
channels are independent, the final state is identical or
$\dist$-equivalent, so both chains realise $\Desc$ and terminate within
$\floor$ of $\BState$. They are distinct chains.
\end{proof}

\begin{theorem}[Path opacity]
\label{thm:opacity}
Two isomorphism chains $\Chain$ and $\Chain'$ realising the same
description $\Desc$ from the same source $\src$ are indistinguishable by
any test of the output state $x_n$.
\end{theorem}

\begin{proof}
Both chains terminate in $\BState_\Desc$ (within $\floor$). Any
test of the output measures properties of $x_n$; by the behaviour cell
definition, all states within $\floor$ of $\BState_\Desc$ exhibit the
same described behaviour (they are members of a region, not a point). The
intermediate states are not in the output and cannot be inspected. Hence
no output test distinguishes the chains.
\end{proof}

Theorem~\ref{thm:opacity} is the formal licence for compiler freedom:
given a behavioural description, the compiler may choose any chain among
those realising it. The user cannot detect, and need not care about, which
chain was chosen.

% ======================================================================
\section{The Covering Floor and Chain Existence}
\label{sec:floor}
% ======================================================================

\begin{definition}[Covering floor]
\label{def:floor}
The \emph{covering floor} of a finite primitive vocabulary $\Phys$ over
state space $(\State,\dist,\Sig)$ is
\[
  \floor \;=\; \inf_{\BState}\;\sup_{x\in\State}\;
  \inf_{(\Eff_1,\dots,\Eff_n)}\;\dist\!\bigl(x_n,\BState\bigr),
\]
where the outer infimum is over all behaviour cells, the supremum over
all source states, and the inner infimum over all finite chains. The floor
$\floor>0$ when $\Phys$ is finite and the state space has more
distinguishable configurations than $\Phys$ can reach in finite steps.
\end{definition}

\begin{theorem}[Floor theorem for effect chains]
\label{thm:floor}
For any finite primitive vocabulary $\Phys$ and any behaviour cell
$\BState$, no finite isomorphism chain reaches the centre of $\BState$
exactly. There exists $\floor>0$ such that for all chains $\Chain$ of
finite length, $\dist(x_n, \BState)\ge\floor$.
\end{theorem}

\begin{proof}
The set of states reachable by a finite chain from $\src$ via $|\Phys|$
primitives is a finite union of images of continuous maps on $\State$.
The centre of $\BState$ (if $\BState$ has nonempty interior) is a single
point; a finite union of images cannot cover it exactly when the
primitives have bounded displacement and the state space has measure
greater than the sum of the displacements. Hence the infimum distance is
positive. The uniform lower bound $\floor>0$ follows from compactness of
$\State$ (or its finite discretisation).
\end{proof}

\begin{corollary}[Chain existence within the floor]
\label{cor:existence}
For every source $\src$, behaviour cell $\BState$, and description
$\Desc\subseteq\Phys$, there exists a finite isomorphism chain of length
at most $|\Desc|\cdot\lceil\Sig/\min_p\delta_p\rceil$ that realises
$\Desc$ and satisfies $\dist(x_n,\BState)\le\floor$.
\end{corollary}

\begin{proof}
Apply each primitive in $\Desc$ in any order, with parameters chosen
greedily to reduce $\dist(\cdot,\BState)$ at each step. Since each
primitive reduces the distance by at least $\delta_p/2$ on the relevant
channel (by continuity of $\dist$ and non-degeneracy of $\BState$), and
the total distance is bounded by $\Sig$, the process terminates within the
claimed bound.
\end{proof}

Corollary~\ref{cor:existence} is the existence guarantee that underlies
the compiler: for any behavioural description expressible in the primitive
vocabulary, a chain exists and the compiler can find it. The floor
$\floor$ is the irreducible residue---the gap between ``exhibits the
behaviour'' and ``is perfectly the behaviour''---which the Floor Theorem
proves cannot be closed by any finite chain.

\begin{remark}[The floor as a feature]
\label{rem:floor-feature}
The covering floor is not a defect of the system. A user watching the
output of a ``water'' chain cannot perceive whether the result is at the
cell centre or at the floor boundary: by the behaviour cell definition,
all states within $\floor$ of $\BState_{\text{water}}$ exhibit the water
behaviour indistinguishably. The floor is perceptually invisible. What it
prevents is not a real failure of the rendering but a category error---the
demand that a finite chain reach a point rather than a cell.
\end{remark}

% ======================================================================
\section{The Catalytic Power of Effects}
\label{sec:power}
% ======================================================================

\subsection{Power of a single effect}

Each effect in a chain closes a fraction of the remaining distance to the
behaviour cell. We formalise this as catalytic power.

\begin{definition}[Catalytic power]
\label{def:power}
The \emph{catalytic power} of an effect $\Eff$ toward cell $\BState$ at
state $x$ is
\[
  \pot(\Eff,x,\BState)
  \;=\;\frac{\dist(x,\BState)-\dist(\Eff(x),\BState)}
            {\dist(x,\BState)-\floor}
  \;\in\;[0,1],
\]
the fraction of above-floor distance to $\BState$ closed by one
application of $\Eff$. An effect is \emph{inert} ($\pot=0$) if it does
not reduce $\dist(\cdot,\BState)$; an effect is an \emph{absolute
catalyst} ($\pot=1$) if it closes all above-floor distance in one step.
\end{definition}

\begin{lemma}[Power is in $[0,1]$]
\label{lem:powerbound}
For any effect, state, and cell, $\pot\in[0,1]$.
\end{lemma}

\begin{proof}
The numerator is at most the denominator because, by the Floor Theorem,
$\dist(\Eff(x),\BState)\ge\floor$, so the distance closed never exceeds
the available above-floor distance. Non-negativity holds for effects
directed toward $\BState$ (those directed away have $\pot=0$ by
convention).
\end{proof}

\subsection{Multiplicative composition}

\begin{theorem}[Multiplicative law]
\label{thm:multiplicative}
For a chain $\Chain=(\Eff_1,\dots,\Eff_n)$ with individual powers
$\pot_1,\dots,\pot_n$, the composite power is
\[
  \pot(\Chain)
  \;=\;1-\prod_{i=1}^{n}(1-\pot_i).
\]
\end{theorem}

\begin{proof}
Let $r_0=\dist(\src,\BState)-\floor$ be the initial above-floor distance.
After $\Eff_1$, the residual above-floor distance is $r_1=(1-\pot_1)r_0$.
After $\Eff_2$, it is $r_2=(1-\pot_2)r_1=(1-\pot_2)(1-\pot_1)r_0$. By
induction, $r_n=\prod_i(1-\pot_i)\,r_0$. The fraction of above-floor
distance closed by the chain is $1-r_n/r_0=1-\prod_i(1-\pot_i)$.
\end{proof}

\begin{corollary}[Diminishing returns]
\label{cor:diminishing}
Repeating a single effect of power $\pot$ a total of $n$ times yields
composite power $1-(1-\pot)^n$, which is strictly increasing in $n$ and
converges to $1$ but never reaches it for finite $n$. Each additional
application closes a strictly smaller absolute fraction of remaining
distance than the previous.
\end{corollary}

The diminishing-returns law is the formal warrant for the design rule:
\emph{diversify effects rather than repeat them}. Adding a new primitive
that addresses a different perceptual channel always adds more composite
power than repeating a primitive already in the chain.

\subsection{Chain saturation}

\begin{theorem}[Saturation dichotomy]
\label{thm:saturation}
A chain applying effects with powers $(\pot_i)_{i\ge1}$ drives the
residual above-floor distance to $0$ if and only if
$\sum_{i=1}^{\infty}\pot_i=\infty$.
\end{theorem}

\begin{proof}
The residual fraction after $n$ effects is $\prod_{i\le n}(1-\pot_i)$.
Taking logarithms, $\sum_{i\le n}\log(1-\pot_i)$ diverges to $-\infty$
(residual $\to0$) if and only if $\sum_i\pot_i=\infty$, since
$-\log(1-\pot)\asymp\pot$ for small $\pot$ and the comparison test
applies. This is a Borel--Cantelli-type criterion
\citep{borel1909,cantelli1917}.
\end{proof}

Theorem~\ref{thm:saturation} distinguishes chains that can in principle
approach the cell arbitrarily closely (non-summable powers: a supply of
effects each addressing a fresh perceptual dimension) from those that
cannot (summable powers: a chain that exhausts available channels and
then adds ever-weaker variants). For the compiler, it implies a
saturation check: given the description $\Desc$, can the chain of
primitives that realises it achieve non-summable composite power? If not,
the chain is statically flagged as insufficient.

% ======================================================================
\section{Coherence of a Chain}
\label{sec:coherence}
% ======================================================================

The multiplicative law assumes that all effects in the chain are directed
toward the same behaviour cell. When effects are directed toward different
cells---or when their combined perceptual character is internally
contradictory---the chain is incoherent and the output is undefined.

\begin{definition}[Support and the support graph]
\label{def:support}
For effects $\Eff_i,\Eff_j$ in a chain and target cell $\BState$, say
$\Eff_j$ \emph{supports} $\Eff_i$ at threshold $\theta\in(0,1]$ if the
perceptual dimension addressed by $\Eff_j$ reinforces, at degree
$\ge\theta$, the shift toward $\BState$ induced by $\Eff_i$. The
\emph{support graph} $G_\Chain$ has an edge $j\to i$ when $\Eff_j$
supports $\Eff_i$ at $\theta$.
\end{definition}

\begin{theorem}[Linear support fails]
\label{thm:linear}
No linear chain of effects in which each effect's contribution is
justified only by the next, and the last is justified by nothing internal
to the chain, can drive residual distance below the floor.
\end{theorem}

\begin{proof}
The terminal effect has no support; its contribution rests on nothing
within the chain. The chain's coherence is therefore contingent on the
terminal effect, which by the Floor Theorem cannot itself resolve to the
cell centre. A linear chain's justification is therefore open at one end,
and by the floor theorem, cannot be closed by any finite extension.
\end{proof}

Theorem~\ref{thm:linear} is the formal content of the following
observation: a chain of effects that are merely \emph{listed}---first
this, then that, then another---with no mutual reinforcement, produces a
result whose behaviour character is contingent on the last effect applied.
Remove the last effect and the character disappears; there is no internal
triangle holding it in place.

\begin{theorem}[Coherence requires a triangle]
\label{thm:triangle}
\leavevmode
\begin{enumerate}[label=(\roman*),leftmargin=1.8em,itemsep=2pt]
\item \emph{(Necessity)} If a chain brings the output within any margin
  of the floor across the behaviour cell, its support graph contains a
  directed cycle of length $\ge3$.
\item \emph{(Sufficiency)} Three effects mutually supporting one another
  at threshold $\theta>\frac{1}{2}$ in a strongly connected triangle are
  coherent and their combined shift toward $\BState$ is robust to removal
  of any single effect.
\end{enumerate}
\end{theorem}

\begin{proof}
\emph{(i)} By Theorem~\ref{thm:linear}, no acyclic support structure
grounds the chain; an acyclic support graph has a source effect with no
incoming support. A $1$-cycle is self-support (vacuous: an effect that
cites only itself). A $2$-cycle is mutual definition between two effects
without an independent third check, which fails the majority condition
$\theta>\frac{1}{2}$ because two effects cannot outvote their own
disagreement. Hence the shortest grounding cycle has length $\ge3$.
\emph{(ii)} In a strongly connected triangle, each effect is supported by
the other two; with $\theta>\frac{1}{2}$, two supporters jointly exceed
any single dissenting pull. Removing one effect leaves the remaining two
still mutually supporting above a reduced but positive threshold.
\end{proof}

\begin{corollary}[Ordinal decidability of coherence]
\label{cor:decide}
Chain coherence is decidable from the signs of pairwise support relations
alone---whether each pair of effects reinforces or undercuts each other's
shift toward $\BState$---with no knowledge of parameter magnitudes or
absolute power values.
\end{corollary}

Corollary~\ref{cor:decide} is the basis for the type-checker described in
\S\ref{sec:compiler}: the assessor can determine whether a chain is
coherent from the relational structure of its effects, before any
parameter is instantiated and before any frame is rendered.

% ======================================================================
\section{The Media Effect Encoder Language}
\label{sec:language}
% ======================================================================

\subsection{Design principles}

The language design follows from the theory:

\begin{enumerate}[leftmargin=1.4em,itemsep=2pt]
\item \textbf{Behaviour, not parameters.} The user writes behavioural
  descriptions; the compiler produces parameter settings. The user's
  vocabulary is physical phenomena, not numerical values.
\item \textbf{Effects as chain steps.} A scene is a \emph{pipeline}: a
  left-to-right sequence of effect steps, each a physical primitive
  description, that the compiler assembles into an isomorphism chain. The
  pipeline is the user's notation for a chain.
\item \textbf{Goal blocks.} The intended behaviour is stated in a
  \texttt{goal} block as a set of behaviour descriptions. The type-checker
  verifies, before rendering, that the pipeline's chain is coherent and
  directed toward the goal.
\item \textbf{Brand invariants.} Named constraints---brand requirements,
  creative guidelines---are declared as \texttt{brand} blocks and enter
  the chain as catalysts with declared confidence weights.
\item \textbf{Multi-format dispatch.} A single source scene can be
  dispatched to multiple output formats (aspect ratios, durations) via a
  \texttt{dispatch} table.
\end{enumerate}

\subsection{Concrete syntax}

A MEE program has the following top-level structure:

\begin{lstlisting}[caption={MEE program structure},label={lst:structure}]
scene NAME {
  clip(PATH, at=T, for=D)
    |> acts_like(DESCRIPTION)
    |> compose(EFFECT, EFFECT, ...)
    |> render(FORMAT)

  goal {
    behaviour: DESCRIPTION
    coherence: >= 0.5
    duration:  < 30s
  }

  brand RULE_NAME {
    invariant: "CONSTRAINT"
    confidence: 0.8
  }

  dispatch {
    when 16:9 do render(widescreen)
    when 9:16  do render(vertical)
    when 1:1   do render(square)
  }
}
\end{lstlisting}

The \texttt{clip} keyword introduces the source footage and anchors the
chain to a physical file, a time offset, and a duration. The pipeline
operator \texttt{|>} composes effects left-to-right, each step taking the
output of the previous as its input. \texttt{acts\_like} introduces a
behaviour description; \texttt{compose} introduces explicit effect
primitives when the user wishes to specify them directly.

\begin{example}[Water surface]
\label{ex:water}
A shot to be transformed so that the scene appears to be viewed through a
water surface:

\begin{lstlisting}[caption={Water surface scene}]
scene water_reveal {
  clip("footage/product.mp4", at=0s, for=15s)
    |> acts_like("water surface")
    |> brand(product_palette, confidence=0.85)
    |> render(motion)

  goal {
    behaviour: "water surface"
    coherence: >= 0.5
    duration:  < 20s
  }

  brand product_palette {
    invariant: "brand blue dominant"
    confidence: 0.85
  }
}
\end{lstlisting}

The compiler resolves \texttt{acts\_like("water surface")} to a
description $\Desc_{\text{water}}=\{\text{oscillate},
\text{refract},\text{reflect},\text{ripple}\}$ and derives an isomorphism
chain realising all four phenomena. The \texttt{brand} catalyst enters the
chain as an effect with the declared confidence weight, contributing to
the composite power multiplicatively.
\end{example}

\begin{example}[Drum skin]
\label{ex:drum}
A shot in which the visual field behaves like a struck drum skin:

\begin{lstlisting}[caption={Drum skin scene}]
scene drum_impact {
  clip("footage/impact.mp4", at=2s, for=8s)
    |> acts_like("drum skin")
    |> compose(
         oscillate(axis=radial, decay=0.3),
         propagate(origin=center, speed=0.8)
       )
    |> render(motion)

  goal {
    behaviour: "drum skin"
    coherence: >= 0.5
  }
}
\end{lstlisting}

Here the user supplements the behavioural description with two explicit
primitive specifications. The compiler merges these into the derived chain,
treating the explicit effects as constraints on the instantiation of
the corresponding phenomena (radial oscillation with a specified decay
rate, radial propagation from the image centre at a specified speed) and
deriving parameter settings for the remaining phenomena in
$\Desc_{\text{drum}}=\{\text{oscillate-radial},\text{propagate-radial},
\text{damp-overtones},\text{membrane-tension}\}$.
\end{example}

\subsection{Behaviour descriptions and the primitive vocabulary}

A behavioural description resolves to a finite set of physical phenomena.
The primitive vocabulary organises phenomena into four namespaces, each
addressing a distinct perceptual channel:

\begin{enumerate}[leftmargin=1.4em,itemsep=2pt]
\item \textbf{Spatial}: transformations of pixel positions.
  Examples: \textit{oscillate} (periodic displacement),
  \textit{distort} (aperiodic warp), \textit{shear} (linear mapping),
  \textit{compress}, \textit{stretch}, \textit{rotate}.
\item \textbf{Photometric}: transformations of colour and light.
  Examples: \textit{refract} (index-of-refraction mapping),
  \textit{scatter} (diffuse light spread), \textit{reflect} (mirror
  composite), \textit{grade} (tone mapping).
\item \textbf{Temporal}: transformations of how the clip evolves over
  time. Examples: \textit{propagate} (wave front advancing in time),
  \textit{damp} (amplitude decay with time), \textit{pulse}
  (periodic temporal envelope), \textit{stutter}.
\item \textbf{Acoustic}: transformations of the audio channel.
  Examples: \textit{resonate} (resonant frequency emphasis),
  \textit{attenuate} (high-frequency cut), \textit{reverb}, \textit{beat}.
\end{enumerate}

A behavioural description such as ``drum skin'' resolves to phenomena
drawn across these namespaces: spatial oscillation (radial), temporal
propagation and damping, and acoustic resonance. The compiler draws on all
four namespaces to build the chain.

\begin{principle}[Behavioural completeness]
\label{prin:completeness}
A chain is \emph{behaviourally complete} for a description $\Desc$ if
every namespace represented in $\Desc$ contributes at least one effect to
the chain. Behavioural completeness is a necessary condition for
Corollary~\ref{cor:existence} (chain existence within the floor).
\end{principle}

% ======================================================================
\section{The Intermediate Representation and Compiler}
\label{sec:compiler}
% ======================================================================

\subsection{The intermediate representation}

\begin{definition}[MEE intermediate representation]
\label{def:ir}
The \emph{MEE intermediate representation} (IR) is an expression language
with four constructors:
\begin{enumerate}[label=(\roman*),leftmargin=1.8em,itemsep=1pt]
\item $\mathsf{Clip}(f, t, d)$: source footage at path $f$, time offset
  $t$, duration $d$.
\item $\mathsf{Prim}(p,\theta)$: a physical primitive $p\in\Phys$
  instantiated at parameter assignment $\theta$.
\item $\mathsf{Compose}(\xi_1,\dots,\xi_n)$: sequential composition of
  $n$ IR expressions, threading output as input.
\item $\mathsf{Hole}(h)$: an unresolved behavioural description $h$,
  type-checked away before any rendering.
\end{enumerate}
A \emph{ground IR expression} contains no \texttt{Hole} constructors. The
compiler's output is always a ground IR expression.
\end{definition}

The IR has three key properties. First, it is compositional: replacing any
subexpression with an evaluation-equal one preserves the semantics of the
whole (the compositionality lemma \citep{scott1970,plotkin1977}). Second,
the presence of \texttt{Hole} constructors is an error at render time---
the type-checker ensures all holes are resolved before the backend is
invoked. Third, the ground IR is substrate-neutral: it can be lowered to
Remotion (TypeScript), to a JSON timeline, or to any other rendering
backend by a pure function from IR to backend syntax.

\subsection{Compiler pipeline}

The compiler is a four-stage pipeline:

\begin{enumerate}[leftmargin=1.4em,itemsep=2pt]
\item \textbf{Tokenise.} A hand-written scanner produces a flat token
  sequence. Keywords, identifiers, operators, string literals, time
  literals, and the pipe operator \texttt{|>} are recognised.
\item \textbf{Parse.} A recursive-descent parser with precedence climbing
  for expressions \citep{pratt1973} produces an AST. Error recovery
  records parse errors and skips to the next recovery keyword, so a
  single bad statement does not abort compilation.
\item \textbf{Type-check and resolve.} The type-checker performs six
  checks and transforms \texttt{Hole} nodes into \texttt{Compose} nodes:
  \begin{enumerate}[label=(\alph*),leftmargin=1.6em,itemsep=1pt]
  \item \emph{Description resolution}: each \texttt{acts\_like} string is
    resolved to a set of physical phenomena $\Desc\subseteq\Phys$; an
    unknown description is an error.
  \item \emph{Coherence check}: the support graph of the resolved chain
    is inspected for a directed cycle of length $\ge3$
    (Theorem~\ref{thm:triangle}). An acyclic or $\le2$-cycle chain emits
    a \texttt{CoherenceWarning} with a suggestion.
  \item \emph{Saturation check}: the chain's composite power is estimated
    from the description's namespace coverage; if coverage is insufficient
    for non-summable power, a \texttt{SaturationWarning} is emitted.
  \item \emph{Brand constraint check}: declared \texttt{brand} invariants
    are resolved against the chain; an invariant whose corresponding
    phenomenon is absent from the chain is an error.
  \item \emph{Goal reachability}: given the chain and the goal block, the
    type-checker predicts whether the chain can bring the source state
    within $\floor$ of the goal cell. If not, it emits a
    \texttt{GoalUnreachable} error with the missing namespaces.
  \item \emph{Duration budget}: the declared duration is checked against
    the sum of clip durations and effect durations; an overrun emits a
    \texttt{DurationExceeded} error.
  \end{enumerate}
\item \textbf{Emit.} A pure function lowers the ground IR to the target
  backend. For Remotion, the emitter produces a \texttt{<Composition>}
  with nested \texttt{<Sequence>} elements, \texttt{interpolate()} calls
  for continuous parameters, and layered components for each perceptual
  channel.
\end{enumerate}

\begin{proposition}[Type-checker soundness]
\label{prop:soundness}
A ground IR expression that passes all six type-checking conditions has a
well-defined denotation: a finite isomorphism chain that is behaviourally
complete for the stated description, coherent (contains a support triangle
of length $\ge3$), and reaches within $\floor$ of the goal cell.
\end{proposition}

\begin{proof}
Condition (a) ensures $\Desc\subseteq\Phys$, so the chain is defined over
real primitives. Condition (b) ensures Theorem~\ref{thm:triangle}(ii) is
satisfied. Conditions (c)--(e) ensure the saturation bound of
Corollary~\ref{cor:existence} and the goal-reachability
Corollary~\ref{cor:existence}. Condition (f) ensures the clip fits within
the declared duration. Together these are the hypotheses of
Theorem~\ref{thm:triangle}(ii) and Corollary~\ref{cor:existence}; the
conclusion follows.
\end{proof}

\subsection{The two-phase architecture}

Following the mutual-exclusion principle between compilation and execution
\citep{turing1936,mccarthy1960}, MEE separates the process into two
strictly ordered phases:

\begin{enumerate}[leftmargin=1.4em,itemsep=2pt]
\item \textbf{Compile.} The tokeniser, parser, type-checker, and emitter
  run to completion, producing a ground IR expression and a backend
  source file. No rendering occurs. The effect registry is fixed and
  immutable at the end of this phase.
\item \textbf{Render.} The backend source file is handed to the rendering
  engine (Remotion). No new effects can be introduced; no new
  descriptions can be resolved. The registry is read-only.
\end{enumerate}

The separation is enforced by the type system: \texttt{Hole} nodes are
representation-level errors at render time. A backend that encounters a
\texttt{Hole} must refuse to render. This ensures that every render is of
a well-typed, coherence-checked, goal-verified composition.

\begin{remark}[Composition inflation]
\label{rem:inflation}
With $d$ perceptual channels and $n$ temporal steps, the number of
distinguishable labelled chains is $d\cdot(1+d)^{n-1}$
\citep{kolmogorov1965}. For $d=4$ (spatial, photometric, temporal,
acoustic) and $n=6$ steps, this gives $4\cdot5^5=12500$ distinct chains.
The expressive space of even a short MEE pipeline is therefore large, and
the type-checker's role---filtering to coherent, goal-directed chains---is
essential to tractable compilation.
\end{remark}

\subsection{The plugin architecture}

New physical primitives are added to the system without modifying the
compiler core. The plugin architecture exposes four stable abstractions:

\begin{enumerate}[leftmargin=1.4em,itemsep=2pt]
\item \textbf{A typed AST.} The four IR constructors
  (\texttt{Clip/Prim/Compose/Hole}) are fixed. Plugins extend the set of
  available $p\in\Phys$ but do not add constructors.
\item \textbf{A primitive trait.} Each $p\in\Phys$ must implement:
  a \emph{signature} (input type, output type, parameter schema), a
  \emph{provider} (the rendering function mapping parameters to backend
  instructions), and a \emph{namespace} (spatial/photometric/temporal/
  acoustic).
\item \textbf{A behaviour registry.} The mapping from description strings
  to $\Desc\subseteq\Phys$ is a registry. New descriptions are registered
  without modifying existing ones.
\item \textbf{An operation registry.} The type-checker resolves each
  primitive by name against the operation registry. Additions are
  well-formedness-preserving.
\end{enumerate}

The frozen nature of the four abstractions is forced by the No-Privileged-
Level property of the underlying framework: a system that treats its own
extension mechanism as a special case---one that can modify the AST---is
self-inconsistent under the same compositional calculus it applies to
user programs. The plugin surface is therefore fixed by theorem, not by
convention.

% ======================================================================
\section{Backward Trajectory Completion as the Compilation Strategy}
\label{sec:backward}
% ======================================================================

The compiler's chain-derivation strategy follows from the backward
trajectory completion insight: the target cell is known at compile time
(from the goal block and behaviour description), and the compiler derives
the isomorphism chain by working backward from the terminus.

\begin{principle}[Backward chain derivation]
\label{prin:backward}
Given a goal cell $\BState_\Desc$ and a source state $\src$:
\begin{enumerate}[label=(\roman*),leftmargin=1.8em,itemsep=1pt]
\item Select the namespace whose current contribution to $\dist(\src,\BState)$
  is largest; choose the primitive from that namespace that maximally
  reduces it.
\item Append the chosen effect to the \emph{front} of the chain-under-
  construction (backward completion: the chain is built from terminus
  toward source).
\item Repeat until all phenomena in $\Desc$ have been addressed and the
  chain's estimated composite power is sufficient for the saturation
  condition.
\end{enumerate}
\end{principle}

The backward construction has two advantages over forward enumeration.
First, it never generates chains that are incoherent from the start:
beginning from the target cell ensures that every prefix of the chain is
already directed toward $\BState$. Second, it is efficient: the number of
steps is bounded by $|\Desc|$ (one phenomenon per step) plus a
saturation tail, rather than by the combinatorial space of all forward
paths.

\begin{remark}[Virtual sub-states]
\label{rem:virtual}
Principle~\ref{prin:backward} permits the compiler to introduce
\emph{virtual sub-states}: intermediate states that are not physically
reachable from $\src$ in a single step from the current position but that
lie on the backward trajectory from $\BState_\Desc$ to $\src$. By the
path opacity theorem (Theorem~\ref{thm:opacity}), these intermediate states
are not observable in the output; they are internal to the compilation
process. The compiler may pass through a state that looks nothing like
either the source or the target, provided the global output lies within
$\floor$ of $\BState_\Desc$. This is the formal licence for locally
``wrong'' intermediate effects---an oscillation step that momentarily
makes the clip look worse before it is combined with a reflection step to
produce the water character.
\end{remark}

The walking analogy (from \S\ref{sec:intro}) is precise here. A person
does not know the sequence of muscle activations required to walk to a
destination; they know the destination and the current position, and the
locomotion system generates the intermediate steps. They notice only when
the trajectory goes wrong---``walking weirdly''---not the intermediate
states. In MEE, the compiler is the locomotion system; the user specifies
destination and source; the type-checker is what notices when the
trajectory goes wrong.

% ======================================================================
\section{Multi-Format Output and the Representation Equivalence}
\label{sec:format}
% ======================================================================

A single MEE scene may be dispatched to multiple output formats. Different
aspect ratios, durations, and platform-specific constraints require
distinct Remotion compositions; the MEE IR is lowered separately for each.

\begin{principle}[One IR, many backends]
\label{prin:backends}
The MEE IR is independent of the rendering backend. A scene compiled once
to a ground IR expression can be emitted as:
\begin{itemize}[leftmargin=1.4em,itemsep=1pt]
\item a Remotion TypeScript composition (widescreen, vertical, square);
\item a JSON timeline (for preview or external tooling);
\item a parameter record (for audit or A/B comparison).
\end{itemize}
The IR expression is the single representation from which all backends
derive.
\end{principle}

This is the representation-equivalence principle applied to rendering:
the same behaviour description can be carried by any backend that can
render the physical primitives in $\Phys$. The choice of backend is
computational, not semantic: it does not change what the scene \emph{does}
to the viewer, only the format in which the instructions are expressed.

The dispatch mechanism in the MEE language (\texttt{dispatch when FORMAT
do render(TARGET)}) is the compile-time routing: for each declared format,
the emitter is invoked once, producing a separate output file. The
backing IR expression is shared; only the emitter-level parameter
adjustments (resampled resolution, re-timed durations, cropped aspect
ratio) differ.

% ======================================================================
\section{Implementation Architecture}
\label{sec:impl}
% ======================================================================

The MEE system is implemented in two tiers:

\paragraph{Web playground (TypeScript).}
A self-contained TypeScript implementation of the full compiler pipeline
---tokeniser, parser, type-checker, and Remotion emitter---runs in the
browser. It exposes a code editor (Monaco), a diagnostic panel (showing
type-checker output), and a compiled-output panel (showing the generated
Remotion composition). The playground is limited to the behaviour
descriptions registered in the browser-side registry; it does not execute
renders. Its purpose is to demonstrate the language surface and the
type-checker's static analysis.

\paragraph{Production compiler (Rust).}
The production compiler is implemented in Rust. It implements the same
four-stage pipeline but with:
\begin{itemize}[leftmargin=1.4em,itemsep=1pt]
\item a richer primitive registry (the full set of Remotion-accessible
  effects and third-party effect libraries);
\item a chain-derivation engine implementing
  Principle~\ref{prin:backward};
\item parallel emission for multi-format dispatch;
\item a JSON result record carrying the IR, the type-checker log, and
  the per-format output paths.
\end{itemize}

The two tiers share the IR definition and the type-checker logic; the
Rust implementation is the normative reference for correctness, and the
TypeScript playground is a projection of it restricted to the browser
environment.

\paragraph{Backend.}
The Remotion backend generates a TypeScript file containing a single
exported \texttt{<Composition>} component. Each \texttt{Compose} node in
the IR becomes a \texttt{<Sequence>} in the Remotion tree; each
\texttt{Prim} node becomes an \texttt{interpolate()} call or a layered
component according to the primitive's namespace. The mapping from IR
nodes to Remotion constructs is a pure function; no state is shared
between the compiler and the renderer.

% ======================================================================
\section{Discussion}
\label{sec:discussion}
% ======================================================================

\paragraph{The user model.}
MEE is designed for a user who knows what they want a shot to \emph{do}
but does not know what parameters achieve it. This is the normal condition
for a video producer: they have seen water, they know how it moves, they
can describe it; they have never set a displacement map's frequency and
damping parameters. The language meets the user at the level of their
knowledge---physical phenomena, behavioural character---and resolves the
gap between description and parameter mechanically.

\paragraph{The assessor as type-checker.}
The type-checker is not a post-hoc lint tool; it is the core of the
system. Every chain that reaches the render phase has already been
certified coherent, behaviourally complete, and goal-directed. The
assessor and the producer are not separate systems; the assessor is
statically embedded in the compilation pipeline. This follows from the
theory: coherence is decidable from ordinal data alone
(Corollary~\ref{cor:decide}), so the check is cheap and can run at
compile time without knowledge of parameter magnitudes.

\paragraph{What the language does not do.}
MEE does not generate footage; it processes existing footage. It does not
determine whether a given shot is \emph{good} in an aesthetic sense; it
certifies that a chain is \emph{coherent} in a structural sense. It does
not guarantee that a rendered output will achieve any perceptual goal in
any specific viewer; by the covering-floor theorem, no finite chain can
guarantee that. What it guarantees is that the chain is well-formed and
directed: if it fails to produce the intended character, the failure is
in the description, not in the chain that realises it.

\paragraph{Ordinal throughout.}
The theory and the language are ordinal. Power values are used internally
by the compiler to decide chain construction strategy and to estimate
saturation; they are never exposed to the user and never asserted as
absolute magnitudes. The user's feedback loop---``more,'' ``less,''
``wrong direction''---is sign-based, matching the ordinal detectability of
Corollary~\ref{cor:decide}. This is not a limitation of the current
implementation; it is a consequence of the floor theorem, which proves
that absolute parameter specifications are unavailable to any finite
observer.

% ======================================================================
\section{Conclusion}
\label{sec:conclusion}
% ======================================================================

A video effect chain is an isomorphism: a finite sequence of physical
primitives connecting a source footage state to a target behaviour cell.
The user specifies behaviour---``acts like water,'' ``acts like a drum
skin''---not parameters. The compiler derives the chain by backward
trajectory completion from the target cell; the type-checker certifies
coherence, saturation sufficiency, and goal reachability before any frame
is rendered; the emitter lowers the chain to Remotion.

The calculus underlying this design is ordinal and structural. Effects
compose multiplicatively (the power law); a chain saturates its target
cell if and only if its aggregate power is non-summable (the Borel--Cantelli
dichotomy); a chain is coherent if and only if its support graph contains
a directed cycle of length at least three. Each of these properties is
checkable from the chain's structural description alone, with no access to
parameter values and no knowledge of what the viewer will perceive. The
compiler is therefore both necessary and sufficient: necessary because the
user cannot specify parameter values, and sufficient because the
structural checks are decidable at compile time.

``Acts like water'' is not a vague instruction. It is a description that
resolves to a finite set of physical phenomena, each with a namespace, a
parameter schema, and a rendering provider. The chain that realises it is
not unique---many chains exhibit water behaviour---but any of them is
correct, and the type-checker certifies correctness before rendering
begins.

\bigskip
\noindent\textbf{Acknowledgements.} This working paper develops a formal
basis for the Media Effect Encoder DSL. It is self-contained and cites
only externally published literature.

\bibliographystyle{plainnat}
\bibliography{references}

\end{document}
